\subsection{אסטרטגיית האימון}

בפרויקט זה לא אומנה מערכת אחת מקצה לקצה, אלא שני מודלי
\textenglish{YOLO} נפרדים, שכל אחד מהם פותר תת־בעיה אחרת:

\begin{itemize}
\item \textbf{מודל זיהוי הרובוטים:} מקבל פריים מלא מן המשחק ומחזיר
  תיבות תיחום לכל הרובוטים יחד עם סיווג לברית כחולה או אדומה.
\item \textbf{מודל זיהוי הספרות:} מקבל חיתוך של רובוט בודד ומחזיר את
  הספרות המופיעות על הבאמפר.
\end{itemize}

שני המודלים אומנו כמודלי \textenglish{Object Detection}. לכן במהלך
האימון המודל למד גם למקם תיבת תיחום סביב האובייקט וגם לסווג את
המחלקה. שלבי \textenglish{ByteTrack}, \textenglish{Levenshtein}
וההצבעה המצטברת אינם מאומנים באמצעות \textenglish{loss}, אלא פועלים
לאחר המודלים כ־\textenglish{post-processing} לוגי.

\subsection{סביבת האימון והיפר־פרמטרים}

אימון שני המודלים בוצע על גבי \textenglish{Google Colab} תוך שימוש
במעבד גרפי \textenglish{NVIDIA T4}, באמצעות ספריית
\textenglish{Ultralytics}. בחירה זו אפשרה להריץ כמה ניסויים בזמן סביר,
לשמור גרפי אימון, ולהשוות בין ריצות שונות בלי להזדקק לחומרה מקומית
חזקה.

ההיפר־פרמטרים העיקריים של ריצות האימון הסופיות מרוכזים להלן:

\begin{itemize}
\item \textbf{\textenglish{Optimizer}:} \textbf{AdamW}.
\item \textbf{\textenglish{Batch Size}:} מודל הרובוטים ---
  \textbf{48}, מודל הספרות --- \textbf{64}.
\item \textbf{מספר \textenglish{Epochs}:} מודל הרובוטים ---
  \textbf{20}, מודל הספרות --- \textbf{150}.
\item \textbf{גודל קלט (\textenglish{Image Size}):} מודל הרובוטים ---
  \textbf{640}, מודל הספרות --- \textbf{128}.
\item \textbf{\textenglish{Patience (Early Stopping)}:} מודל הספרות --- \textbf{30}.
\item \textbf{קצב למידה התחלתי:} מודל הרובוטים --- \textbf{0.001667},
  מודל הספרות --- \textbf{0.000714}. (נבחר אוטומטית על ידי
  \textenglish{YOLO})
\item \textbf{\textenglish{מומנטום}:} \textbf{0.9}. (נבחר אוטומטית על ידי
  \textenglish{YOLO})
\end{itemize}

גודל הקלט השונה בין שני המודלים נובע מאופי המשימה. מודל הרובוטים מקבל
פריים מלא ולכן אומן בגודל 640, בעוד שמודל הספרות מקבל חיתוך קטן של
אזור הבאמפר לאחר זיהוי הרובוט. גודל 128 נבחר משום שהוא קרוב לגודל
הממוצע של החיתוכים שהתקבלו בפועל ולגודל הצפוי בזמן הרצה על סרטוני
תחרות אמיתיים. למרות הקלט הקטן יותר, מודל הספרות אומן למשך יותר אפוקים
משום שמשימת זיהוי הספרות קשה יותר וכוללת עשר מחלקות, אובייקטים קטנים
וחיתוכים לא מושלמים.

\subsection{אימון מודל זיהוי הרובוטים}

מודל זיהוי הרובוטים אומן על מערך נתונים קיים מ־\textenglish{Roboflow}
הכולל מעל 1000 תמונות, עם שתי מחלקות בלבד: ברית כחולה וברית אדומה.
מבחינת למידת מכונה, זה היה המודל הפשוט יותר מבין השניים, משום שהאובייקט
שצריך לזהות הוא יחסית גדול, והרמז החזותי של צבע הברית בולט יותר לעין מן
הספרות הקטנות שעל הבאמפר.

מטרת האימון הייתה להגיע לשתי יכולות בסיסיות ואמינות:

\begin{itemize}
\item \textbf{לוקליזציה טובה:} תיבות תיחום יציבות ומדויקות של כל
  הרובוטים בפריים.
\item \textbf{סיווג ברית אמין:} הבחנה נכונה בין רובוט כחול לרובוט אדום.
\end{itemize}

בזכות פשטות יחסית זו, המודל התכנס מהר יחסית והגיע לתוצאה גבוהה של
\textenglish{mAP@50} של כ־\textbf{0.97}. תוצאה זו הספיקה כדי להפוך את
המודל לבסיס יציב עבור שלב המעקב: גם אם הזיהוי אינו מושלם בכל פריים,
הוא מדויק מספיק כדי לספק ל־\textenglish{ByteTrack} קלט איכותי.

האתגרים המרכזיים במודל זה היו בעיקר תנאי צילום משתנים, חסימות חלקיות
בין רובוטים וטשטוש תנועה בזמן האצה. כדי לשפר את העמידות למצבים אלו,
נעשה שימוש ב־\textenglish{Augmentation} של הטיה, שינוי בהירות והוספת
רעש, כפי שתואר בפרק מערך הנתונים.

\subsection{אימון מודל זיהוי הספרות}

מודל זיהוי הספרות היה המודל המאתגר יותר בפרויקט. בניגוד למודל
הרובוטים, כאן יש עשר מחלקות שונות (0--9), והאובייקטים קטנים בהרבה,
לעיתים מטושטשים, חתוכים חלקית או מצולמים בזווית חדה. בנוסף, הספרות
מופיעות על רובוטים בתנועה, ולעיתים החיתוך של הרובוט עצמו אינו מושלם.

מסיבות אלו תהליך האימון של מודל הספרות היה תהליך איטרטיבי יותר, שבו
המשימה המרכזית הייתה לא רק לשפר את המודל אלא גם לשפר את איכות
הנתונים. חלק מן התמונות הראשונות נאספו ידנית ממקורות איכותיים יחסית,
אך בהמשך נוספו גם חיתוכים אוטומטיים מסרטוני משחק אמיתיים. תוספת זו הייתה
חשובה במיוחד, משום שהיא קירבה את מערך האימון להתפלגות הנתונים שהמודל
פוגש בפועל בזמן הרצה.

למרות הקושי, המודל הגיע ל־\textenglish{Precision} של כ־\textbf{0.75}
ברמת הספרה הבודדת.  זהו ביצוע נמוך יותר מזה של מודל הרובוטים, אך הוא
עדיין שימושי במערכת המלאה, משום שזיהוי הספרות אינו מקור ההחלטה היחיד:
בהמשך הוא משולב עם צבע הברית, לוח המשחקים והצבעה מצטברת לאורך זמן.

גם כאן בוצעה \textenglish{Augmentation}, אך היא הייתה אגרסיבית וממוקדת
יותר: חיתוך חלקי, סיבוב, שינוי בהירות, שינוי חשיפה והטיה. אוגמנטציות
אלו נבחרו משום שהן מדמות היטב את תנאי ההפעלה בפועל, שבהם הבאמפר אינו
מופיע תמיד בדיוק באותה זווית או באותו מקום בתמונה.

\subsection{תוצאות האימון}

באיורים הבאים מופיעים גרפי האימון של המודלים, המציגים את השינוי במדדים
ובפונקציות ה־\textenglish{Loss} בסטי האימון והוולידציה לאורך האפוקים:

\begin{figure}[H] \centering
  \includegraphics[width=0.8\textwidth]{robots-results}
  \caption{גרפי האימון של מודל זיהוי הרובוטים.}
\end{figure}

\begin{figure}[H] \centering
  \includegraphics[width=0.8\textwidth]{digits-results}
  \caption{גרפי האימון של מודל זיהוי הספרות.}
\end{figure}

בטבלה הבאה מופיעות תוצאות האימון עבור כל אחד מהמודלים, שנמדדו על סט
הבדיקה שלא שימש לאימון או כוונון המודלים:

\begin{table}[H]
  \centering
  \begin{tabular}{|c|c|c|p{0.4\textwidth}|} \hline
    \textbf{מדד} & \textbf{מודל הרובוטים} & \textbf{מודל הספרות} & \textbf{משמעות המדד} \\ \hline
    Precision & 0.93 & 0.75 & מתוך כלל הזיהויים שהמודל החזיר, כמה היו נכונים. \\ \hline
    Recall & 0.94 & 0.78 & מתוך כלל האובייקטים האמיתיים, כמה המודל הצליח לזהות. \\ \hline
    mAP@50 & 0.97 & 0.85 & איכות הזיהוי כאשר חיזוי נחשב נכון אם החפיפה עם התיבה האמיתית היא לפחות 0.5 \textenglish{IoU}. \\ \hline
    mAP@50-95 & 0.66 & 0.43 & מדד מחמיר יותר שבודק איכות זיהוי בטווח ספי \textenglish{IoU} שונים. \\ \hline
  \end{tabular}
\end{table}

שני המודלים אומנו כמודלי \textenglish{YOLO} ולכן נמדדו באמצעות אותם
מדדי זיהוי אובייקטים. עם זאת, יש לפרש את המדדים בהקשר שונה: מודל
הרובוטים בוחן זיהוי של רובוטים וצבע ברית, ואילו מודל הספרות בוחן זיהוי
של ספרות קטנות על גבי רובוטים, משימה רגישה יותר לטשטוש, חסימות וזוויות
צילום.

מטריצות הבלבול (\textenglish{Confusion Matrix}) של המודלים ותוצאות
אימון מלאות לפי מחלקה מופיעות בנספח \ref{app:training}.

גרפי האימון, טבלת התוצאות ומטריצות הבלבול מצביעים על הבדל ברור בין שני
המודלים. מודל זיהוי הרובוטים יציב וחזק יותר, משום שההבחנה בין ברית
אדומה לכחולה מבוססת על סימן חזותי גדול וברור יחסית ברוב
הפריימים. לעומת זאת, מודל הספרות מתמודד עם משימה קשה יותר: הספרות
קטנות, מופיעות על באמפרים של רובוטים, ולעיתים מצולמות בזווית, בטשטוש
תנועה, בהסתרה חלקית או בתאורה לא אחידה. לכן גם כאשר מודל הספרות מצליח
לזהות ספרות רבות, יש יותר מקום לבלבול בין מחלקות, והמערכת אינה מסתמכת
על חיזוי מפריים יחיד אלא משלבת את תוצאות המודל לאורך זמן באמצעות
\textenglish{tracking}, צמצום לפי לוח המשחקים והצבעה מצטברת.

\subsection{ניסויים ושיפורים לאורך הדרך}

במהלך הפיתוח בוצעו כמה ניסויים ושינויים מהותיים. גם כאשר לא נשמרו כל
המספרים המדויקים של כל ריצה, אפשר לזהות בבירור את כיוון השיפור בכל שלב:

\begin{itemize}
\item \textbf{ניסוי 1 --- אימון ראשוני על נתונים נוחים וזיהוי בעיית
    \textenglish{Variance}:} בתחילת הדרך, חלק גדול מן הדוגמאות של
  הספרות הגיע ממקורות איכותיים יחסית. בשלב זה זיהיתי תופעה של
  \textenglish{High \newline Variance (Overfitting)}: המודל למד לזהות
  היטב את תמונות האימון והגיע לדיוק גבוה, אך כשהורץ על נתוני ולידציה
  המורכבים מסרטוני תחרות אמיתיים (עם טשטוש תנועה, חסימות ותאורה חלשה),
  הביצועים צנחו משמעותית. המודל למד להתאים את עצמו באופן הדוק מדי
  לתנאים האידיאליים ולא הצליח להכליל (\textenglish{Generalize}) את
  יכולת הזיהוי לתנאי אמת. העובדה שהמודל הגיע לביצועים גבוהים על נתוני
  האימון מעידה שלא סבלנו מבעיית \textenglish{High Bias
    (Underfitting)}.

\item \textbf{ניסוי 2 --- הוספת חיתוכים אוטומטיים כאמצעי
    \textenglish{Regularization}:} כדי לצמצם את הפער בין שגיאת האימון
  לשגיאת הוולידציה, נוספו למערך הנתונים דוגמאות שנחתכו אוטומטית מתוך
  סרטוני משחק. מערך הנתונים נעשה "מלוכלך" וריאלי יותר. הוספת רעש
  ויזואלי לנתוני האימון שימשה הלכה למעשה כסוג של רגולריזציה
  (\textenglish{Regularization}): היא מנעה מהמודל להסתמך רק על
  מאפיינים מושלמים, הכריחה אותו ללמוד מאפיינים עמידים
  (\textenglish{Robust Features}), וכך הקטינה את
  ה־\textenglish{Variance} ושיפרה את ההתאמה בין שלב האימון לזמן ההפעלה
  בפועל.
\item \textbf{ניסוי 3 --- כוונון אוגמנטציה:} אוגמנטציה מתונה מדי לא
  הספיקה כדי להתמודד עם זוויות קשות ותאורה משתנה, אך אוגמנטציה חזקה מדי
  גרמה לעיתים לפגיעה בקריאות הספרות. לכן נדרש איזון בין גיוון הנתונים
  לבין שמירה על צורת הספרה.
\item \textbf{ניסוי 4 --- פתרון מערכתי במקום שיפור מודל בלבד:} בשלב
  מסוים התברר שגם אם מודל הספרות לא מגיע לדיוק מושלם, עדיין אפשר לשפר
  משמעותית את המערכת כולה באמצעות שילוב התוצאה עם צבע הברית, לוח המקצים
  והצבעה מצטברת לאורך זמן. זהו שינוי חשוב בגישת הפיתוח: מעבר משיפור
  מודל בודד לשיפור מערכת מלאה.
\end{itemize}

\subsection{אתגרים מרכזיים ופתרונות}

תהליך האימון חשף כמה אתגרים מרכזיים:

\begin{itemize}
\item \textbf{פער בין נתוני אימון לנתוני אמת:} תמונות נקיות או חיתוכים
  ידניים לא תמיד מייצגות נאמנה את סרטוני התחרות. הפתרון היה להוסיף יותר
  דוגמאות מתוך צינור העבודה האמיתי.
\item \textbf{אובייקטים קטנים ודומים:} הספרות קטנות, וחלקן דומות זו
  לזו במיוחד כאשר יש טשטוש תנועה. הפתרון היה גם הרחבת הנתונים וגם
  הסתמכות על לוגיקת החלטה מערכתית לאחר האימון.
\item \textbf{סיכון ל־\textenglish{Overfitting}:} במיוחד במודל הספרות,
  גודל מערך הנתונים היה מוגבל יחסית. שימוש באוגמנטציה ובדוגמאות מגוונות
  יותר צמצם את הבעיה.
\item \textbf{מגבלת משאבים:} סביבת \textenglish{Colab} נוחה לאימון,
  אך עדיין דורשת תכנון זהיר של מספר הניסויים ושל זמני הריצה. לכן חלק מן
  האיטרציות בוצעו תחילה על קבוצות נתונים קטנות יותר או עם בדיקות מהירות
  לפני ריצות מלאות.
\end{itemize}

\subsection{זמן אימון וחומרה}

כל ריצות האימון בוצעו על גבי \textenglish{GPU} מסוג
\textenglish{NVIDIA T4} ב־\textenglish{Google Colab}.

\begin{itemize}
\item \textbf{זמן אימון סופי של מודל הרובוטים:} \textbf{20.59 דקות}.
\item \textbf{זמן אימון סופי של מודל הספרות:} \textbf{10.92 דקות}.
\item \textbf{מספר ריצות ניסוי כולל:} \textbf{כ־8--10}.
\end{itemize}

גם ללא המספרים המדויקים, חשוב לציין שהחלק הארוך באמת בפרויקט לא היה רק
זמן האימון עצמו, אלא זמן האיסוף, התיוג והשיפור האיטרטיבי של הנתונים.
במובן זה, הפרויקט היה \textbf{Data-Centric} לא פחות משהיה
\textbf{Model-Centric}.

\subsection{סיכום תהליך האימון}

תהליך האימון הוביל לשתי מסקנות עיקריות. ראשית, מודל זיהוי הרובוטים היה
חזק יחסית כבר בשלבים מוקדמים והגיע במהירות לביצועים מספקים. שנית,
האתגר האמיתי בפרויקט היה מודל זיהוי הספרות, שבו איכות הנתונים והריאליזם
של הדוגמאות היו קריטיים לא פחות מהארכיטקטורה עצמה.

לכן, ההצלחה של המערכת אינה נובעת רק מבחירת מודל טוב, אלא משילוב בין
אימון סביר של שני מודלים, שיפור מתמיד של מערך הנתונים, ושכבת החלטה
נוספת שמנצלת את הפלטים של המודלים בצורה חכמה. זהו בדיוק המעבר מפרויקט
של "אימון רשת" לפרויקט של הנדסת מערכת מבוססת למידת מכונה.
